% TOP_PROBLEM: {"id":30007532,"problem_number":"LOCAL-30007532","title":"State-dependent fiscal multipliers","category_id":21,"category":{"id":21,"name":"economics","display_name":"Economics","description":"Open economic theory statements, published research agendas, and proposed scoped specifications; question provenance and historical priority are qualified per record.","slug":"economics","order_index":21,"created_at":"2026-10-08T00:00:00Z"},"status":"research_question","external_url":null,"legacy_ids":[],"published":false,"statement":"Identify cumulative fiscal multipliers conditional on economic states and shock size under comparable financing and monetary-policy rules.","economics":{"original_id":21,"field":"Fiscal policy and sovereign debt","kind":"identification","formalization_basis":"adapted_research_question","source_status":"research_agenda","priority_status":"earliest_found","priority_note":"This is the earliest explicit relevant statement verified in this audit, not a claim of original priority; older literature and earlier versions may contain antecedents.","earliest_verified_reference":{"authors":["Adrien Auclert","Matthew Rognlie","Ludwig Straub"],"title":"Fiscal and Monetary Policy with Heterogeneous Agents","year":2025,"url":"https://web.stanford.edu/~aauclert/annual_review_hank.pdf","doi":"10.1146/annurev-economics-091624-044646","locator":"Section 5, pp. 19–21 of the April 2025 author manuscript; “Nonlinearities” and “Optimal policy”","evidence_summary":"The review says empirical size and sign asymmetries remain debated. Its optimal-policy discussion identifies unresolved computation and equity–efficiency issues and warns that unrestricted monetary–fiscal Ramsey problems with standard preferences need not have a well-defined steady state."},"current_reference":{"authors":["Adrien Auclert","Matthew Rognlie","Ludwig Straub"],"title":"Fiscal and Monetary Policy with Heterogeneous Agents","year":2025,"url":"https://web.stanford.edu/~aauclert/annual_review_hank.pdf","doi":"10.1146/annurev-economics-091624-044646","locator":"Section 5, pp. 19–21 of the April 2025 author manuscript; “Nonlinearities” and “Optimal policy”","evidence_summary":"The review says empirical size and sign asymmetries remain debated. Its optimal-policy discussion identifies unresolved computation and equity–efficiency issues and warns that unrestricted monetary–fiscal Ramsey problems with standard preferences need not have a well-defined steady state."},"formalization_note":"The review explicitly notes unresolved empirical nonlinearities and size/sign asymmetries. The richer state-conditioned multiplier is an adapted, not canonical, specification."},"statusQualification":"Published question or research agenda verified at the cited locator. The mathematical specification is adapted for this catalogue and includes additional assumptions; source publication does not establish first-ever priority or certify this exact model remains unsolved. The review explicitly notes unresolved empirical nonlinearities and size/sign asymmetries. The richer state-conditioned multiplier is an adapted, not canonical, specification.","statusReviewedAt":"2026-10-08"}
% FIRST_POSTED: 2026-10-08
% ENTRY_KIND: statement_only
% !TeX program = lualatex
\documentclass[11pt]{article}
\usepackage{fontspec}
\usepackage[a4paper,margin=25mm]{geometry}
\usepackage{amsmath,amssymb,mathtools}
\usepackage{xurl}
\usepackage[hidelinks]{hyperref}
\newcommand{\sref}[2]{\hyperlink{source:#1:#2}{[S#2]}}
\setlength{\emergencystretch}{3em}
\begin{document}
% problemId:30007532
\section*{State-dependent fiscal multipliers}\label{30007532}
\subsection{Definitions and mathematical statement}
Let \(G_t\) be real government purchases and \(Y_t\) real output. Define predetermined state \(s_t\) using a predeclared recession indicator, inflation, debt, openness, and monetary-policy constraints. An instrument \(Z_t\) identifies an unanticipated purchase innovation, excluding responses to anticipated future output and other simultaneous fiscal measures. For intervention magnitude \(a\), let \(Y_{t+h}(a)\) and \(G_{t+h}(a)\) be potential outcomes under a stated financing and monetary-policy rule. Define the cumulative state-dependent multiplier
\[
m_H(s,a)=
\frac{\sum_{h=0}^H E[Y_{t+h}(a)-Y_{t+h}(0)\mid s_t=s]}
{\sum_{h=0}^H E[G_{t+h}(a)-G_{t+h}(0)\mid s_t=s]},
\]
when its denominator is nonzero. Identify the response surface for a fixed horizon and observed range of shock sizes, with confidence sets robust to weak instruments and overlapping outcomes. Determine which state differences remain after holding instrument composition, financing, and policy reactions comparable. Distinguish marginal from large-intervention multipliers and positive from negative spending shocks; linear local projections cannot automatically establish those equalities. Report support limitations when large shocks occur only in exceptional crises. The goal is a transportable conditional estimate under explicit identification assumptions, not one invariant fiscal multiplier applicable to every economy and spending program.

\par\textbf{Formulation and scope.} The review explicitly notes unresolved empirical nonlinearities and size/sign asymmetries. The richer state-conditioned multiplier is an adapted, not canonical, specification.

\par\textbf{Open-question provenance.} An explicit question or research agenda was verified at the source locator. This does not by itself certify that every later version remains unresolved \sref{30007532}{1}.

\par\textbf{Historical priority.} This is the earliest explicit relevant statement verified in this audit, not a claim of original priority; older literature and earlier versions may contain antecedents.

\subsection{Short English statement}
Identify cumulative fiscal multipliers conditional on economic states and shock size under comparable financing and monetary-policy rules.

\subsection{Sources}
\begin{itemize}\raggedright
\item[\textnormal{[S1]}]\hypertarget{source:30007532:1}{} Adrien Auclert, Matthew Rognlie, Ludwig Straub. 2025. Fiscal and Monetary Policy with Heterogeneous Agents. Section 5, pp. 19–21 of the April 2025 author manuscript; “Nonlinearities” and “Optimal policy”.
\url{https://web.stanford.edu/~aauclert/annual_review_hank.pdf}.
DOI: 10.1146/annurev-economics-091624-044646.
{\small\textit{Earliest verified question/agenda reference; historical priority qualified below. The review says empirical size and sign asymmetries remain debated. Its optimal-policy discussion identifies unresolved computation and equity–efficiency issues and warns that unrestricted monetary–fiscal Ramsey problems with standard preferences need not have a well-defined steady state.}}
\end{itemize}
\end{document}
